\section{SEC-014: Secrets in Git History via Memory Tool}
\label{sec:sec-014}

\subsection*{Metadata}
\begin{tabular}{ll}
\textbf{Issue ID} & SEC-014 \\
\textbf{Severity} & High (CVSS 8.2) \\
\textbf{Category} & Security \\
\textbf{CWE} & CWE-200: Exposure of Sensitive Information \\
\textbf{Affected File} & \srcfile{src/tools/memory.ts}, \srcfile{src/tools/sandbox-memory.ts} \\
\textbf{Branch} & \branchref{fix/SEC-014-secrets-git-history} \\
\textbf{Changes} & \branchdiff{fix/SEC-014-secrets-git-history} \\
\textbf{Commit} & 8667321 \\
\end{tabular}

\subsection*{Executive Summary}
Every \texttt{save} operation in the \texttt{memory} tool writes content to a file and commits it to git via \texttt{git add \&\& git commit}. If the saved content contains sensitive information---API keys, tokens, passwords, private data---that information becomes permanently embedded in the git repository's history with no automated mechanism to remove it. Even if a subsequent save removes the secret from the current content, the old commit still contains it. Standard git garbage collection (\texttt{git gc}) does not remove historical commits containing secrets.

Unlike SEC-002 (which focuses on detection of secrets before they are committed), SEC-014 addresses the permanence of secrets once they reach git history. The fix introduces automatic git history squashing: when the number of memory commits exceeds a threshold (50), the entire memory file history is rewritten into a single commit using \texttt{git filter-branch}, effectively removing all historical content. Additionally, archive files are no longer git-tracked, and lightweight \texttt{git gc} is run after each save to prune unreachable objects.

\subsection*{Root Cause Analysis}
Each memory save creates a git commit:

\begin{lstlisting}[language=TypeScript]
// memory.ts:153-160
await mkdir(dirname(memoryFile), { recursive: true });
await writeFile(memoryFile, finalContent, "utf-8");

try {
    execSync(`git add "${memoryPath}" && git commit -m "..."`, {
        cwd, stdio: "pipe",
    });
} catch (err: any) {
\end{lstlisting}

The content is committed as a blob, tree, and commit object in git. These objects are immutable---once created, they persist until explicitly garbage-collected. An attacker with repository access can search all git history:

\begin{lstlisting}[language=bash]
git log -p --all -S "sk-proj" -- memory/MEMORY.md
\end{lstlisting}

The \texttt{archiveOverflow()} function also git-added archive files, creating additional permanent records of older content:

\begin{lstlisting}[language=TypeScript]
try {
    execSync(`git add "${archiveFile}"`, { cwd, stdio: "pipe" });
} catch {
    // Not in git, that's fine
}
\end{lstlisting}

\subsection*{Fix Implementation}
The fix introduces three protections. First, a commit-counting mechanism triggers history squashing when the threshold is exceeded:

\begin{lstlisting}[language=TypeScript]
const MAX_MEMORY_COMMITS = 50;

function countMemoryCommits(cwd: string, memoryPath: string): number {
    try {
        const out = execSync(
            \texttt{git log --oneline -- "\${memoryPath}" 2>/dev/null | wc -l},
            { cwd, encoding: "utf-8", stdio: "pipe" },
        ).trim();
        return parseInt(out, 10) || 0;
    } catch { return 0; }
}

async function squashMemoryHistory(
    cwd: string, memoryPath: string, memoryFile: string,
    commitCount: number,
): Promise<void> {
    const content = await readFile(memoryFile, "utf-8");
    const branch = getCurrentBranch(cwd);

    try {
        execSync(
            \texttt{git filter-branch --force --prune-empty --index-filter } +
            \texttt{'git rm --cached --ignore-unmatch "\${memoryPath}"' HEAD},
            { cwd, stdio: "pipe", timeout: 60000, maxBuffer: 10 * 1024 * 1024 },
        );
        execSync(\texttt{git update-ref -d refs/original/refs/heads/\${branch} } +
            \texttt{2>/dev/null; rm -rf .git/refs/original}, { cwd, stdio: "pipe" });
        await writeFile(memoryFile, content, "utf-8");
        execSync(\texttt{git add "\${memoryPath}" && git commit -m "Memory history } +
            \texttt{squashed (\${commitCount} commits \textrightarrow{} 1)"},
            { cwd, stdio: "pipe" });
        execSync(\texttt{git gc --aggressive --prune=now},
            { cwd, stdio: "pipe", timeout: 120000 });
    } catch (err: any) {
        console.error(\texttt{SEC-014: Memory history squash failed: \${err.message}});
    }
}
\end{lstlisting}

Second, archive files are no longer git-added:

\begin{lstlisting}[language=TypeScript]
// SEC-014: Do NOT git-add archive files --- they contain old content that
// may include secrets. Archive stays on disk but is not version-controlled.
\end{lstlisting}

Third, the integration point in \texttt{execute()} checks the commit count after each save:

\begin{lstlisting}[language=TypeScript]
const commitCount = countMemoryCommits(cwd, memoryPath);
if (commitCount >= MAX_MEMORY_COMMITS) {
    await squashMemoryHistory(cwd, memoryPath, memoryFile, commitCount);
}
await runGitGC(cwd);
\end{lstlisting}

The same archive-add removal and \texttt{git gc} call are applied in \texttt{sandbox-memory.ts}.

\subsection*{Affected Files}
\begin{itemize}
    \item \srcfile{src/tools/memory.ts} --- Added \texttt{MAX\_MEMORY\_COMMITS} threshold (50), \texttt{countMemoryCommits()} counting function, \texttt{squashMemoryHistory()} filter-branch based squashing, \texttt{runGitGC()} for periodic pruning, removed git-add from \texttt{archiveOverflow()}, integrated squash check in \texttt{execute()} after each save.
    \item \srcfile{src/tools/sandbox-memory.ts} --- Removed git-add from \texttt{archiveOverflow()}, added \texttt{git gc --auto --quiet} after saves.
    \item \srcfile{test/memory-sec-014.test.ts} --- New standalone test suite for git history protections.
\end{itemize}

\subsection*{Verification}
Standalone test suite validates all git history protection mechanisms: commit counting, history squashing, content preservation after squash, garbage collection execution, archive file exclusion from git tracking, and reflog cleanup.

\subsection*{Test File}
\srcfile{test/memory-sec-014.test.ts}

Contains 7 test cases: \texttt{countMemoryCommits}, squash reduces commits, content preserved, gc, archive files not git-added, reflog destroyed.
